\makeatletter
\newcommand{\setword}[2]{%
  \phantomsection
  #1\def\@currentlabel{\unexpanded{#1}}\label{#2}%
}
\makeatother

\newcommand{\boolNo}{\text{No}} % boolean No
\newcommand{\boolYes}{\text{Yes}} % boolean Yes

\NewDocumentCommand{\localG}{g}{%
    \IfNoValueTF{#1}
    {G}%
    {G_{#1}}%
}
\NewDocumentCommand{\sparseG}{g}{%
    \IfNoValueTF{#1}
    {\widetilde{G}}%
    {\widetilde{G}_{#1}}%
} % sparsified graph in dynamic agreement (DA)
\NewDocumentCommand{\dyCC}{g}{%
    \IfNoValueTF{#1}
    {\textsc{CC}}%
    {\textsc{CC}_{#1}}%
} % dynamic data structure for connected components
\NewDocumentCommand{\journal}{g}{%
    \IfNoValueTF{#1}
    {\mathcal{J}}%
    {\mathcal{J}_{#1}}%
}
\newcommand{\epsDA}{\varepsilon_{\rm DA}} % accuracy parameter for dynamic agreement
\NewDocumentCommand{\interestSet}{g}{%
    \IfNoValueTF{#1}
    {\mathcal{I}}%
    {\mathcal{I}_{#1}}%
} % interesting set in DA
\NewDocumentCommand{\anchorSet}{g}{%
    \IfNoValueTF{#1}
    {\Phi}%
    {\Phi_{#1}}%
} % anchor set in DA


\newcommand{\openN}[2]{N_{#1}(#2)} % open neighborhood of vertex #2 in graph #1
\newcommand{\closedN}[2]{N_{#1}[#2]} % open neighborhood of vertex #2 in graph #1
\newcommand{\snapshotN}[1]{\widetilde{N}(#1)} % open neighborhood of vertex #2 in graph #1


\newcommand{\explicitDA}{\mathrm{ExplicitDynamicAgreement}} % explicit dynamic agreement algorithm
\newcommand{\ADA}{\mathrm{AmplifiedDynamicAgreement}} % explicit dynamic agreement algorithm
\newcommand{\DA}{\mathrm{DynamicAgreement}} % explicit dynamic agreement algorithm
\NewDocumentCommand{\Notify}{g}{%
    \IfNoValueTF{#1}
    {\mathrm{Notify}}%
    {\mathrm{Notify}(#1)}%
}
\NewDocumentCommand{\Anchor}{g}{%
    \IfNoValueTF{#1}
    {\mathrm{Anchor}}%
    {\mathrm{Anchor}(#1)}%
}
\NewDocumentCommand{\Clean}{g}{%
    \IfNoValueTF{#1}
    {\mathrm{Clean}}%
    {\mathrm{Clean}(#1)}%
}
\NewDocumentCommand{\Connect}{g}{%
    \IfNoValueTF{#1}
    {\mathrm{Connect}}%
    {\mathrm{Connect}(#1)}%
}
\NewDocumentCommand{\LabelSync}{g}{%
    \IfNoValueTF{#1}
    {\mathrm{LabelSync}}%
    {\mathrm{LabelSync}(#1)}%
}
\NewDocumentCommand{\ProbAgree}{g}{%
    \IfNoValueTF{#1}
    {\mathrm{ProbablisticAgreement}}%
    {\mathrm{ProbablisticAgreement}(#1)}%
}
\NewDocumentCommand{\Heavy}{g}{%
    \IfNoValueTF{#1}
    {\mathrm{Heavy}}%
    {\mathrm{Heavy}(#1)}%
}

\newcommand{\delt}[1]{\mathrm{Del}_{#1}}
\newcommand{\addt}[1]{\mathrm{Add}_{#1}}
\newcommand{\delIncEdge}[1]{\delta(#1)} % number of edges incident on vertex #1 that are deleted (\in \delt)
\NewDocumentCommand{\ccclass}{g}{%
    \IfNoValueTF{#1}
    {C}%
    {C_{#1}}%
} % connected component class
\NewDocumentCommand{\ccclasses}{g}{%
    \IfNoValueTF{#1}
    {\mathcal C}%
    {\mathcal{C}_{#1}}%
} % connected component class
\newcommand{\comp}[1]{\mathrm{comp}(#1)} % connected component of vertex #1
\NewDocumentCommand{\piece}{ g }{%
    \IfNoValueTF{#1}
    {piece}%
    {piece#1}%
}
\newcommand{\retainedPiece}[1]{R_{#1}} % retained pieces of class #1
\newcommand{\piecesToMerge}{\mathcal{Z}} %
\newcommand{\pieceToMerge}{P} %
\newcommand{\complabel}{\ell} %

\newcommand{\successDA}{\mathcal{E}_{\rm DA}} % success event of dynamic agreement
\newcommand{\jnTime}[1]{\card{\journal(#1)}} % length of the journal at time #1
\newcommand{\nChangedVers}[1]{c_{#1}} % number of vertices whose labels changed at time #1

\newcommand{\edSolver}{\mathcal D}
\newcommand{\edFactor}{\eta} % approximation factor of the explicit correlation-clustering algorithm
\newcommand{\edTime}[1]{T_{\mathcal D}(#1)}
\newcommand{\edcost}[1]{\widehat{\cost}_{#1}}
\newcommand{\edmean}[2]{\widehat{\cost}^{(#1)}_{#2}}
\newcommand{\edscore}[2]{\widehat f_{#1}(#2)}
\newcommand{\edhistory}[2]{\widehat G_{#1}(#2)}
\newcommand{\edminimum}[1]{\widehat M_{#1}}
\newcommand{\edproxy}[2]{Z_{#1}(#2)}
\newcommand{\edlog}[2]{\mathcal L_{#1}(#2)}
\newcommand{\edbucket}[2]{B_{#1}(#2)}
\newcommand{\edatom}[1]{B_{123}(#1)}

% name for variable used in the implementation of LabelSync
\newcommand{\setClasses}{\mathit{Set}_{\ccclasses}} % variable used in the
\newcommand{\groupdict}{\mathit{group\_dict}} % variable used in the
\newcommand{\compdict}{\mathit{comp\_dict}} % variable used in the
\newcommand{\changedList}{\mathit{changed\_list}} % variable used in the


\section{Explicit Dynamic Algorithm for Insertion-Only Streams}\label{sec:explicit_dynamic_algorithm_for_insertion_only_stream}

In this section, we consider the explicit dynamic model for insertion-only streams.
In this model, at each update time $ t $, a new vertex $ v_{t} $ arrives.
The algorithm needs to maintain the solution and reports $ v_t $ and its cluster label, and for each vertex whose cluster identifier has changed, reports the vertex and its new cluster identifier.

Our results is an explicit dynamic algorithm for consensus clustering in the insertion-only stream model, which achieves a strictly better than $ 2 $ approximation ratio with probablity at least $ 1-\zeta $, and polylogarithmic amortized update time in expectation.
\ExplicitDynamic*

We use the vertex-arrival model of \cref{subsec:online-model}, with the input-label access of \cref{subsec:dynamic-model}.
In particular, the input partitions and arrival order are oblivious, the horizon $ n $ is known, and any selected input label of an arriving vertex can be read in constant time.

We obtain \cref{thm:explicit.dynamic} from the following two lemmas.

\begin{lemma}\label{lem:explicit.consensus.clustering}
Suppose that $ \edSolver $ is a randomized explicit dynamic correlation-clustering algorithm in the database model~(\cref{def:query.model.label}) that, on every oblivious stream of at most $ n $ vertex insertions, has expected amortized update time $ \edTime{n}=\polylog(n+2) $ including queries and reports, and with probability at least $ 1/2 $ maintains an $ \edFactor $-approximation simultaneously at all prefixes.
Then, for every $ 0<\eps<2/(3\edFactor+5) $ and $ \zeta\in(0,1) $, there is an explicit dynamic consensus-clustering algorithm for insertion-only streams that achieves approximation ratio $ 2-2/(3\edFactor+5)+\eps $ at all prefixes simultaneously with probability at least $ 1-\zeta $.
Its expected amortized update time is
\begin{align}
    O\!\left(
        \eps^{-3}\log^2\frac{2n}{\zeta}
        \log\frac{2n}{\eps\zeta}\cdot\polylog(n+2)
    \right).
    \label{eq:explicit.consensus.time}
\end{align}
\end{lemma}

\begin{remark}
Given the upper bound $ \edTime{n} $, we can convert the algorithm of \cref{lem:explicit.consensus.clustering} into one whose amortized update time satisfies \eqref{eq:explicit.consensus.time} on every execution, preserving the approximation guarantee with probability at least $ 1-\zeta $.
\end{remark}

\begin{definition}[Database model~\cite{assadi-wang}]\label{def:query.model.label}
At time $ t $, the algorithm accesses the graph of positive edges on $ V_t $ through three queries.
A degree query returns the number of neighbors of a given vertex, an edge query tests whether two given vertices are adjacent, and a neighborhood-sample query returns an independent uniformly random neighbor of a given vertex of positive degree.
Each query costs $ O(\log(|V_t|+2)) $ time.
\end{definition}

\paragraph{From clique buckets to the database model.}\label{par:explicit.bucket.queries}
By the bucket representation of \cref{subsec:offline-buckets,subsec:dynamic-algorithm}, two distinct vertices of a three-input majority graph are adjacent exactly when their input labels agree on some pair of coordinates.
For a vertex $ v $, the intersection of any two of its three pair buckets is the set of vertices agreeing with $ v $ on all three coordinates.
Order each pair bucket by the remaining coordinate and then by vertex index in an order-statistic tree, so this intersection occupies one rank interval.
The open neighborhood is the disjoint union of the common intersection with $ v $ removed and the complement of the intersection in each of the three buckets.
Summing these four sizes gives the degree, and a uniform rank in their concatenation gives a uniform neighbor.
Thus insertions and all three database queries take $ O(\log(|V_t|+2)) $ time, using three bucket memberships per vertex.

\begin{restatable}{lemma}{ExplicitDA}\label{lem:explicit.correlation.clustering}
For some absolute constant $ \edFactor\ge1 $, there is an explicit randomized dynamic correlation clustering algorithm in the database model~(\cref{def:query.model.label}) achieves an $ \edFactor $-approximation simultaneously at all prefixes of any fixed insertion-only stream with probability at least $ 1-n^{-2} $.
%The algorithm uses $ N $ as its size parameter and takes the majority of three independent runs of each agreement or heaviness test.
For every $ t\le n $, the expected total work through insertion $ t $, including all reported label assignments, is $ O(t\polylog(n+2)) $.
\end{restatable}

\begin{proof}[Proof of~\cref{thm:explicit.dynamic}]
Apply \cref{lem:explicit.consensus.clustering} with the algorithm of \cref{lem:explicit.correlation.clustering}.
For fixed $ \eps $ and inverse-polynomial $ \zeta $, \eqref{eq:explicit.consensus.time} is polylogarithmic in $ n $.
\end{proof}

We defer the proof of \cref{lem:explicit.correlation.clustering} to \cref{sec:explicit_dynamic_correlation_clustering_algorithm_for_insertion_only_stream}.
The remainder of this section is devoted to prove \cref{lem:explicit.consensus.clustering}.
The cases $ n=1 $ and $ m\le2 $ are handled by following an input, as in \cref{sec:online}, so assume $ n\ge2 $ and $ m\ge3 $.

\subsection{Algorithm}\label{subsec:explicit.consensus.algorithm}

We follow the history-score selection rule in \cref{sec:oncon-algorithm}, with the following two modifications.
For every triple candidate, maintain its solution by employing $ \edSolver $ instead of $ \Agree $.
In \cref{sec:oncon-algorithm}, the algorithm decides whether to retain the selected candidate or switch to another candidate using history scores computed from exact costs.
As maintaining the exact cost of each candidate is too expensive in dynamic model, we maintain an estimated cost of each candidate using a set of evaluation samples. 
We detail this estimation below.

We reuse the parameter choices from \cref{sec:oncon-algorithm,subsec:oncon-analysis}:
\begin{equation*}
  c=\eta+1,\qquad
  w=\frac{3c}{3c+2},\qquad
  R_0=2-\frac{2}{3\eta+5},\qquad
  R=R_0+\frac{\eps}{2},
\end{equation*}
with $ \eta $ the approximation factor of $ \edSolver $. 
The candidate set is defined as in \cref{sec:oncon-algorithm}, with the parameters changed by a constant factor as follows.

\begin{align}
    \theta=\gamma=\frac{\eps}{40}, \qquad
    K=\left\lceil\frac{8}{\eps}\log\frac{4n}{\zeta}\right\rceil \label{eq:explicit.dynamic.consensus.K}
\end{align}

Sample $ K $ candidates independently as in \cref{sec:oncon-algorithm}: with probability $ w $, follow a uniform input, and otherwise run $ \edSolver $ on the majority graph of a uniform triple of distinct input indices.
Let $ \A=[K] $, and let $ \A_{\rm in} $ and $ \A_{\rm tri} $ be the sets of indices of input and triple candidates, respectively. As in \cref{sec:oncon-algorithm}, write $ A_a(t) $ for candidate $ a $'s partition and $ \lambda_t^a(v) $ for its local identifier of vertex $ v $.

\paragraph{Evaluation samples.}
Set
\begin{align}
    L=\left\lceil\log_2\frac{4n}{\zeta}\right\rceil, \qquad
    g=\left\lceil\frac{2}{\gamma^2}\log\frac{8KnL}{\zeta}\right\rceil \qquad
    s=Lg.
    \label{eq:explicit.consensus.parameters}
\end{align}

Independently of the candidates, we sample $ L $ groups $ (J_{j,1},\ldots,J_{j,g}) $, with all $ J_{j,b} $ independent and uniform in $ [m] $ as evaluation sets use throughout the course of the algorithm.
For a partition $ C $ of $ V_t $, define
\begin{align}
    \edmean{j}{t}(C)&:=\frac1g\sum_{b=1}^g d(C,P_{J_{j,b}}^t),
    &\edcost{t}(C)&:=\min_{j\in[L]}\edmean{j}{t}(C),\label{eq:explicit.consensus.estimates}\\
    \edscore{a}{t}&:=\edcost{t}(A_a(t)),
    &\edhistory{a}{t}&:=\max_{u\le t}\edscore{a}{u},
    &\edminimum{t}&:=\min_{a\in\A}\edhistory{a}{t}.\nonumber
\end{align}
Initially, $ \edhistory{a}{0}=0 $.

The updating rule of the algorithm are given in \cref{alg:explicit.consensus}.

\tienlong{Unreadable beyond this point. The details are correct though.}

\begin{algorithm}[htbp]
    \DontPrintSemicolon
    \KwIn{The arriving vertex $ v_t $, the sampled indices, and the maintained candidate and selection states.}
    \KwOut{The identifier of $ v_t $ and every changed output identifier.}
    Read and store the labels of $ v_t $ at the sampled input indices\;
    \ForEach{$ a\in\A $}{
        Advance candidate $ a $ to time $ t $; for a triple candidate, first update its majority-graph representation\;
        Let $ \edlog{a}{t} $ contain the assignment to $ v_t $ and the reported changes to old vertex labels\;
        Update all group costs by processing $ \edlog{a}{t} $\;
        Compute $ \edscore{a}{t} $ and set $ \edhistory{a}{t}\gets\max\{\edhistory{a}{t-1},\edscore{a}{t}\} $\;
    }
    Set $ \edminimum{t}\gets\min_{a\in\A}\edhistory{a}{t} $ and apply the selection rule in \cref{sec:oncon-algorithm}, using estimated history scores, to choose $ a_t $\;
    \eIf{$ t=1 $ or $ a_t\ne a_{t-1} $}{
        Assign and report $ (a_t,\lambda_t^{a_t}(v)) $ for every $ v\in V_t $\;
    }{
        Apply and report the assignments in $ \edlog{a_t}{t} $, prefixing each local label by $ a_t $\;
    }
    \caption{Explicit dynamic consensus-clustering algorithm.}
    \label{alg:explicit.consensus}
\end{algorithm}

\subsection{Approximation Guarantee}\label{subsec:explicit.consensus.approximation}

For the analysis, define
\begin{align}
    \edproxy{a}{t}:=
    \begin{cases}
        \cost_t(P_{i_a}^t),&a\in\A_{\rm in},\\
        \OPT_t+c|B_{Q_a}^t(F_t^*)|,&a\in\A_{\rm tri}.
    \end{cases}
    \label{eq:explicit.consensus.proxy}
\end{align}
Call $ a $ good for $ t $ if $ \edproxy{a}{t}\le R\OPT_t $ and, when $ a $ is a triple candidate, its solver achieves the $ \edFactor $-approximation at every prefix.

\begin{lemma}[Existence of a good candidate]\label{lem:explicit.good.candidate}
With probability at least $ 1-\zeta/4 $, every prefix $ t\le n $ has a good candidate.
\end{lemma}
\begin{proof}
The proxy depends only on the sampled candidate type and input indices.
The input expectation in \eqref{eq:oncon-input-history}, the majority-error bound in \eqref{eq:oncon-majority-error}, and the cancellation in the proof of \cref{lem:oncon-expectation} give $ \E[\edproxy{a}{t}]\le R_0\OPT_t $.
The fixed-prefix Markov argument in the proof of \cref{lem:oncon-good-history} therefore gives
\begin{align}
    \Prb[\edproxy{a}{t}\le R\OPT_t]\ge\frac{\eps}{4}.
    \label{eq:explicit.consensus.proxy.success}
\end{align}
If $ \OPT_t=0 $, the proxy is zero for every draw, so this bound still holds.

Conditional on any fixed triple, its solver succeeds throughout the stream with probability at least $ 1/2 $.
Consequently, a candidate is good for $ t $ with probability at least $ \eps/8 $, independently across candidates.
By the choice of $ K $ in \eqref{eq:explicit.dynamic.consensus.K}, the probability that any prefix has no good candidate is at most
\begin{align}
    n\exp(-\eps K/8)\le\zeta/4.
    \label{eq:explicit.consensus.candidate.event}
\end{align}
\end{proof}

\begin{lemma}[Accuracy of the estimated costs]\label{lem:explicit.estimated.costs}
With probability at least $ 1-\zeta/2 $, simultaneously for every $ t\le n $ and $ a\in\A $,
\begin{align}
    \edcost{t}(F_t^*)&\le2\OPT_t,\label{eq:explicit.consensus.ceiling}\\
    \left|
        \edcost{t}(A_a(t))-\edcost{t}(F_t^*)
        -\bigl(\cost_t(A_a(t))-\OPT_t\bigr)
    \right|
    &\le\gamma d(A_a(t),F_t^*).
    \label{eq:explicit.consensus.comparison}
\end{align}
\end{lemma}
\begin{proof}
To prove \eqref{eq:explicit.consensus.ceiling}, apply Markov's inequality to each group mean for the fixed partition $ F_t^* $.
Each mean exceeds $ 2\OPT_t $ with probability at most $ 1/2 $, so their minimum does so with probability at most $ 2^{-L} $.
When $ \OPT_t=0 $, all these means are zero.
A union bound over $ t $ gives failure probability at most $ \zeta/4 $.

For \eqref{eq:explicit.consensus.comparison}, condition on all candidate samples and solver choices, leaving the evaluation samples independent.
For each $ a,t,j $, the difference of the two group means averages independent variables $ d(A_a(t),P_{J_{j,b}}^t)-d(F_t^*,P_{J_{j,b}}^t) $.
By the triangle-inequality argument, each variable lies in $ [-d(A_a(t),F_t^*),d(A_a(t),F_t^*)] $.
Hoeffding's inequality bounds the probability of a deviation larger than $ \gamma d(A_a(t),F_t^*) $ in either direction by $ 2\exp(-g\gamma^2/2) $.
If the distance is zero, the comparison is exact.
A union bound over $ KnL $ comparisons gives failure probability at most $ \zeta/4 $.
Finally, if the differences of two lists of numbers all lie in the same interval, the difference of their minima lies in that interval as well.
Taking the minimum over the groups proves \eqref{eq:explicit.consensus.comparison}.

\end{proof}

\begin{lemma}[Estimated history of a good candidate]\label{lem:explicit.estimated.history}
If $ b $ is good for $ t $ and \eqref{eq:explicit.consensus.comparison} holds at $ t $, then
\begin{align}
    \edminimum{t}\le\edhistory{b}{t}
    \le\edcost{t}(F_t^*)+(R-1+3\gamma)\OPT_t.
    \label{eq:explicit.consensus.good.history}
\end{align}
\end{lemma}
\begin{proof}
The first inequality follows from the definition of $ \edminimum{t} $.
If $ b $ is an input candidate, each group cost is nondecreasing under insertions, so $ \edhistory{b}{t}=\edcost{t}(P_{i_b}^t) $ by the input-cost monotonicity argument in the proof of \cref{lem:oncon-expectation}.
Since $ \cost_t(P_{i_b}^t)\le R\OPT_t $, the triangle inequality gives $ d(P_{i_b}^t,F_t^*)\le(R+1)\OPT_t\le3\OPT_t $.
Substituting these bounds in \eqref{eq:explicit.consensus.comparison} proves the second inequality for an input candidate.

For a successful triple candidate, the derivation of \eqref{eq:oncon-history-distance} uses only the solver's simultaneous approximation guarantee, so it gives $ d(A_b(u),F_t^*|_{V_u})\le c|B_{Q_b}^t(F_t^*)| $ for every $ u\le t $.
Each group mean, and hence their minimum, satisfies both the cost comparison and restriction inequalities used in \cref{lem:oncon-history}.
Applying that proof with estimated costs yields
\begin{align}
    \edhistory{b}{t}
    \le\edcost{t}(F_t^*)+c|B_{Q_b}^t(F_t^*)|
    \le\edcost{t}(F_t^*)+(R-1)\OPT_t.
    \nonumber
\end{align}
\end{proof}

\begin{lemma}[Approximation at every prefix]\label{lem:explicit.selected.cost}
With probability at least $ 1-\zeta $, the output satisfies $ \cost_t(C_t)\le(R_0+\eps)\OPT_t $ simultaneously for every $ t\le n $.
\end{lemma}
\begin{proof}
By \cref{lem:explicit.good.candidate,lem:explicit.estimated.costs} and a union bound, with probability at least $ 1-3\zeta/4 $ every prefix has a good candidate and satisfies \eqref{eq:explicit.consensus.ceiling} and \eqref{eq:explicit.consensus.comparison}.
On this event, fix $ t $ and apply \cref{lem:explicit.estimated.history} to a good candidate at $ t $.
The proof of \cref{lem:oncon-selection} uses only the monotonicity of history scores and the fact that a current score is at most its history score.
It therefore applies without change and gives $ \edcost{t}(C_t)\le(1+\theta)\edminimum{t} $.
Combine this inequality with \eqref{eq:explicit.consensus.good.history}, \eqref{eq:explicit.consensus.comparison} for the selected candidate, and $ d(C_t,F_t^*)\le\cost_t(C_t)+\OPT_t $ to obtain
\begin{align}
    (1-\gamma)\cost_t(C_t)
    &\le(1+\gamma)\OPT_t+\theta\edcost{t}(F_t^*)
        +(1+\theta)(R-1+3\gamma)\OPT_t\nonumber\\
    &\le\bigl(R+4\gamma+\theta(R+1+3\gamma)\bigr)\OPT_t.
    \label{eq:explicit.consensus.final.cost}
\end{align}
The last inequality uses \eqref{eq:explicit.consensus.ceiling}.
Since $ \theta=\gamma=\eps/40<1/160 $ and $ R<2 $,
\begin{align}
    \frac{R+4\gamma+\gamma(R+1+3\gamma)}{1-\gamma}
    &=R+\frac{\gamma(2R+5+3\gamma)}{1-\gamma}
    <R+10\gamma
    =R_0+\frac{3\eps}{4}
    <R_0+\eps.
    \nonumber
\end{align}
This holds for every prefix on the stated event, including prefixes with $ \OPT_t=0 $.
\end{proof}

\subsection{Implementation and Update Time}\label{subsec:explicit.consensus.implementation}


\begin{lemma}[Maintaining the estimated costs]\label{lem:explicit.cost.maintenance}
For each candidate, its $ L $ evaluation-group means can be maintained in $ O(ns) $ words and $ O(s\log(n+2)) $ time per reported label assignment, where $ s=Lg $.
Computing the estimated cost from the updated means takes $ O(L) $ time.
\end{lemma}
\begin{proof}
Apply the counters and update formulas of \cref{subsec:dynamic-algorithm} separately to each evaluation group, with group size $ g $.
Restore the input-only term following \eqref{eq:dynamic-score} to obtain the full group mean before taking its minimum or updating a history score.
That term cannot be omitted here, since it varies across groups and times.
For each candidate, keep its previous labels while processing the arriving vertex's assignment and then the old vertices' reported changes.
A change to a previously unused label is an insertion into an empty block.
As in \cref{subsec:dynamic-algorithm}, each assignment or move takes $ O(s\log(n+2)) $ time, including updates to all group sums.
Each of the $ s $ evaluation positions has at most $ n $ nonempty block-intersection records, giving the stated space bound.
Computing the minimum over the $ L $ groups takes $ O(L) $ time per candidate per arrival.
\end{proof}

\begin{lemma}[Cost of rebuilding the explicit output]\label{lem:explicit.rebuilding}
On every execution, the number $ S $ of outer switches satisfies
\begin{align}
    S\le K\left(2+\left\lceil\log_{1+\theta}(gn^2)\right\rceil\right)
    =O\!\left(K\eps^{-1}\log(gn+2)\right).
    \label{eq:explicit.consensus.switches}
\end{align}
Over the first $ t $ insertions, rebuilding the output at switches takes $ O(t(S+1)) $ time.
\end{lemma}
\begin{proof}
Every estimated score, history score, and minimum history score belongs to $ g^{-1}\mathbb Z\cap[0,\binom n2] $.
Applying the phase-count argument in the proof of \cref{lem:oncon-recourse}, with minimum positive score $ 1/g $ and the present value of $ \theta $, gives \eqref{eq:explicit.consensus.switches}.
Each rebuild writes and reports at most $ t $ identifiers during the first $ t $ insertions.
Between switches, output writes and reports are charged to assignments reported by the selected candidate.
\end{proof}

\begin{lemma}[Total update time]\label{lem:explicit.total.work}
For every $ t\le n $, the expected total time for the first $ t $ insertions, including initialization and all output reports, is
\begin{align}
    O\!\left(
        tK(\edTime{n}+1)s\log(n+2)
        +tK\eps^{-1}\log(gn+2)
    \right).
    \label{eq:explicit.consensus.total.work}
\end{align}
\end{lemma}
\begin{proof}
For every fixed choice of candidate samples, the expected work of all solver copies over the first $ t $ insertions, including their queries and reported assignments, is $ O(tK\edTime{n}) $.
The bucket representation implements each database query in $ O(\log(n+2)) $ time, and \cref{lem:explicit.cost.maintenance} charges $ O(s\log(n+2)) $ time to each reported assignment.
Reading the sampled labels, updating input candidates, and computing scores and selections fit in $ O(tKs\log(n+2)) $ time.
By \cref{lem:explicit.rebuilding}, output rebuilds contribute $ O(tK\eps^{-1}\log(gn+2)) $ time.
Adding these bounds proves \eqref{eq:explicit.consensus.total.work}.
\end{proof}

\begin{lemma}[Bounding the work of each solver copy]\label{lem:explicit.capped.work}
Given the prefix expectation bound in \cref{lem:explicit.consensus.clustering}, one can modify \cref{alg:explicit.consensus} so that \eqref{eq:explicit.consensus.time} bounds the amortized update time at every prefix on every execution.
The modified algorithm retains the simultaneous approximation guarantee of \cref{lem:explicit.selected.cost}.
\end{lemma}
\begin{proof}
Double $ K $ in \eqref{eq:explicit.dynamic.consensus.K} and use this larger value when setting $ g $.
Set $ h=1+\lceil\log_2 n\rceil $.
At insertion $ t $, stop a triple solver permanently if its next operation would bring its cumulative work above $ 8ht\edTime{n} $.
Count only the solver's own work against its budget, and buffer its reports until the insertion finishes.
If the solver stops, discard its pending reports and use the first input partition for that candidate from this insertion onward.
Cache that input's labels as vertices arrive, and rebuild the candidate's labels and score counters once when it switches to that input.
Use fresh identifiers for its input blocks, report the resulting assignments, and retain the candidate's previous history score when taking the running maximum.

Fix a sampled triple and let $ W(u) $ be the cumulative work of its uninterrupted solver through insertion $ u $.
For any $ t\le n $, the first endpoint $ u\ge t $ in $ \{\min(2^j,n):0\le j<h\} $ satisfies $ u\le2t $.
If the budget is exceeded at $ t $, monotonicity gives $ W(u)>4hu\edTime{n} $.
Since $ \E[W(u)]\le u\edTime{n} $, Markov's inequality and a union bound over the at most $ h $ endpoints bound the probability of stopping by $ 1/4 $.
Thus, conditional on any triple, its solver both succeeds at all prefixes and finishes the entire stream with probability at least $ 1/2-1/4=1/4 $.
Using this probability in the proof of \cref{lem:explicit.good.candidate} gives success probability at least $ \eps/16 $ per candidate at each fixed prefix, for which the doubled $ K $ suffices.
The modified trajectories are independent of the evaluation samples, so \cref{lem:explicit.estimated.costs,lem:explicit.estimated.history,lem:explicit.selected.cost} apply as before.

Through insertion $ t $, each solver spends at most $ 8ht\edTime{n} $ operations and rebuilds at most once, at a cost of $ O(ts\log(n+2)) $ for that rebuild.
Consequently the accounting of \cref{lem:explicit.total.work} applies on every execution with $ 8h\edTime{n} $ in place of $ \edTime{n} $.
The additional factor $ h=O(\log(n+2)) $ is absorbed by the polylogarithmic factor in \eqref{eq:explicit.consensus.time}.
\end{proof}

\begin{proof}[Completion of the proof of~\cref{lem:explicit.consensus.clustering}]
The approximation guarantee follows from \cref{lem:explicit.selected.cost}.
By \eqref{eq:explicit.dynamic.consensus.K} and \eqref{eq:explicit.consensus.parameters},
\begin{align}
    K&=O\!\left(\eps^{-1}\log\frac{2n}{\zeta}\right),
    &L&=O\!\left(\log\frac{2n}{\zeta}\right),
    &g&=O\!\left(\eps^{-2}\log\frac{2n}{\eps\zeta}\right).
    \nonumber
\end{align}
Substituting $ s=Lg $ and $ \edTime{n}=\polylog(n+2) $ into \eqref{eq:explicit.consensus.total.work} proves \eqref{eq:explicit.consensus.time}.
\end{proof}

\section{Explicit Dynamic Correlation Clustering Algorithm for Insertion-Only Stream}\label{sec:explicit_dynamic_correlation_clustering_algorithm_for_insertion_only_stream} % (fold)

In this section, we prove the following lemma.
\ExplicitDA*

We present the algorithm $ \explicitDA $, which is an explicit algorithm for correlation clustering.
We demonstrate that $ \explicitDA $ obtains the properties stated in \cref{lem:explicit.correlation.clustering}.

Throughout this section, we denote by $ G $ the underlying graph of the problem, and by $ G_{t} $ the induced subgraph of $ G $ on the first $ t $ vertices.
\subsection{Algorithm \texorpdfstring{$ \explicitDA $}{ExplicitDA}}\label{sub:algorithm.explictDA} % (fold)
The algorithm $ \explicitDA $ (\cref{alg.explicit.dynamic.agreement}) builds on an amplified version of $ \DA $ (\cite{CLMP2024Dynamic}), denoted by $ \ADA $, with additional data structures and procedures for reporting label changes.

$ \DA $ is a randomized dynamic correlation-clustering algorithm.
For each fixed time $ t $, its maintained clustering is an $ \edFactor $-approximation with probability at least $ 1-5/n $, for an absolute constant $ \edFactor $ \cite[Theorem~6]{CLMP2024Dynamic}.
Its expected amortized update time is $ \polylog(n) $ \cite[Appendix~C, proof of Theorem~41]{CLMP2024Dynamic}.
The algorithm does not report the vertices whose cluster labels change after an update; instead, its output procedure enumerates the entire clustering when requested \cite[Appendix~D]{CLMP2024Dynamic}.

We amplify $ \DA $ to obtain $ \ADA $, which maintains an $ \edFactor $-approximation simultaneously at all prefixes with probability at least $ 1-n^{-2} $, while retaining $ \polylog(n) $ expected amortized update time.
The amplification and its analysis are given in \cref{sub:amplify.da}.

Whenever a vertex $ v_t $ arrives, $ \explicitDA $ executes $ \ADA $ to update the sparse graph and anchor set.
To report the resulting label changes, it maintains the following additional data structures:
\begin{itemize}
    \item A journal $ \journal $ records every edge insertion, edge deletion, and change of anchor membership during the execution of $ \ADA $.
    For each edge modified during the update, the journal also records whether it was present immediately before its first modification.

    We denote by $ \journal{t} $ the changes recorded in the journal during the execution of $ \ADA $ on the arrival of $ v_t $.

    \item The dynamic connectivity structure of \cite{holm2001poly}, denoted by $ \dyCC $, maintains the connected components of $ \sparseG $ under edge insertions and deletions.
    It supports each edge update in $ O(\log^2 n) $ amortized time and a query $ \comp{x} $ in $ O(\log n) $ time.
    At a fixed moment, $ \comp{x}=\comp{y} $ if and only if $ x $ and $ y $ belong to the same component of the current graph.
    These component identifiers need not remain stable between updates.

    \item After $ \ADA $ finishes an update, we use a $ \LabelSync $ procedure, which takes as input the journal $ \journal $, the $ \dyCC $ data structure, to maintain a data structure that assigns labels to arrived vertices.
        At each time step $ t $, two vertices $ x $ and $ y $ have the same label if and only if they belong to the same connected component of the current sparse graph $ \sparseG{t} $.
        The procedure $ \LabelSync $ updates the component identifiers of the vertices in $ \sparseG $ to reflect the connected components of the current sparse graph, and reports every vertex whose component identifier changed during update $ t $.
    On the good event defined below, two vertices have the same label exactly when they belong to the same component of $ \sparseG{t} $.
\end{itemize}

\begin{algorithm}[htbp]
    \SetKwInOut{KwIn}{Input}
    \SetKwInOut{KwOut}{Output}
    \KwIn{A known horizon $ n $ and an oblivious stream of at most $ n $ vertex insertions, accessed through the database model of \cref{def:query.model.label}.}
    \KwOut{After each insertion, the new vertex and its label, followed by every previously arrived vertex whose label changed and its new label.}
    Initialize the maintained sparse graph $ \sparseG \gets \emptyset $

    Initialize the dynamic connectivity structure $ \dyCC $

    Initialize $ \LabelSync $

    \On{arrival of $ v_t $}{
        Initialize an empty journal $ \journal{t} $ and insert $ v_t $ into $ \dyCC $ as an isolated vertex

        Execute $ \ADA(v_t) $, recording every edge operation and anchor-membership change in $ \journal{t} $ and applying every edge operation to $ \dyCC $

        Execute $ \LabelSync{t} $ using $ \journal{t} $ and $ \dyCC $
    }
    \caption{$ \explicitDA $: Explicit Dynamic Agreement Algorithm for Insertion-Only Streams}
    \label{alg.explicit.dynamic.agreement}
\end{algorithm}

To prove \cref{lem:explicit.correlation.clustering}, we state the following results about $ \ADA $ and $ \LabelSync $.

\begin{restatable}{lemma}{AmplifiedDynamicAgreementLemma}
    \label{lem:ada.guarantees}
    Fix an oblivious insertion-only stream of at most $ n $ vertices.
    For sufficiently large $ n $, with probability at least $ 1-n^{-2} $, the connected components of the sparse graph maintained by $ \ADA $ give an $ \edFactor $-approximation simultaneously at every prefix.
    For every $ t\le n $, the expected total work of $ \ADA $ through insertion $ t $ is $ O(t\polylog(n)) $.
\end{restatable}

The following lemma summarizes the correctness and implementation cost of $ \LabelSync $, established in \cref{lem:labelsync.classes.are.connected.components,lem:labelsync.cost}.

\begin{restatable}{lemma}{LabelSyncLemma}
    \label{lem:labelsync.guarantees}
    For sufficiently large $ n $ and any fixed oblivious insertion-only stream processed by $ \ADA $, with probability at least $ 1-n^{-2} $, $ \LabelSync $ correctly reports the labels of vertices whose labels changed during each update.
    This holds jointly with the approximation guarantee in \cref{lem:ada.guarantees}.
    For every execution, the work of $ \LabelSync $ during update $ t $, including all reports, is
    \begin{align}
        O((\jnTime{t}+1)\log n+\nChangedVers{t}),
    \end{align}
    where $ \jnTime{t} $ is the journal length and $ \nChangedVers{t} $ counts changed labels of previously arrived vertices.
\end{restatable}

\begin{proof}[Proof of~\cref{lem:explicit.correlation.clustering}]
    By \cref{lem:ada.guarantees}, with probability at least $ 1-n^{-2} $, $ \ADA $ maintains an $ \edFactor $-approximation.
    On each of these executions, by~\cref{lem:labelsync.guarantees}, the labels maintained by $ \LabelSync $ describe exactly the clustering maintained by $ \ADA $ at every prefix.

    Fix a prefix of $ t $ insertions, and let $ W_t $ be the work performed by $ \ADA $ during this prefix.
    \cref{lem:ada.guarantees} gives
    \begin{align}
        \mathbb E[W_t]=O(t\polylog(n+2)).
    \end{align}
    Recording an edge operation or an anchor-membership change takes constant additional work.
    Every journal entry records an operation performed by $ \ADA $, so the total journal length
    \begin{align}
        H_t:=\sum_{s=1}^{t}\jnTime{s}
    \end{align}
    satisfies $ H_t=O(W_t+t) $ on every execution.
    In particular, $ \sum_{s=1}^{t}\card{\delt{s}}\le H_t $, and $ \mathbb E[H_t]=O(t\polylog(n+2)) $.

    By \cref{lem:number.of.updated.vertices}, the total number of changed labels of previously arrived vertices satisfies
    \begin{align}
        \sum_{s=1}^{t}\nChangedVers{s}
        &=O\left(t\log n+\log n\sum_{s=1}^{t}\card{\delt{s}}\right) \\
        &=O((H_t+t)\log n).
    \end{align}
    The newly inserted vertices contribute another $ t $ reported label assignments.
    Summing the work bound in \cref{lem:labelsync.guarantees} over the prefix, the total work of $ \LabelSync $, including its reports, is therefore $ O((H_t+t)\log n) $.
    There are at most $ H_t $ journalled edge operations, each requiring $ O(\log^2 n) $ amortized time in $ \dyCC $ by \cite{holm2001poly}; together with vertex initialization, maintaining $ \dyCC $ costs $ O((H_t+t)\log^2 n) $ time over the prefix.
    Thus the total work of $ \explicitDA $ through insertion $ t $ is
    \begin{align}
        O(W_t+(H_t+t)\log^2 n).
    \end{align}
    Taking expectations yields $ O(t\polylog(n+2)) $, including all reported label assignments.
    The graph queries are implemented in the database model of \cref{def:query.model.label}.
    For the remaining bounded horizons, we can maintain an optimal clustering by exhaustive computation with constant work per insertion, including all reports.
\end{proof}

\subsection{Amplify \texorpdfstring{$ \DA $}{DynamicAgreement} Algorithm}
\label{sub:amplify.da}
We describe $ \ADA $.

We first review $ \DA $.
At every time $ t $, $ \DA $ maintains a sparse subgraph $ \sparseG{t} $ of the input graph $ G_t $ and an anchor set $ \anchorSet{t} $ satisfying
\begin{align}
    \text{every edge of } \sparseG{t} \text{ has an endpoint in } \anchorSet{t}. \tag{I} \label{eq:invariant.anchor}
\end{align}
The maintained clustering is the partition into connected components of $ \sparseG{t} $.

Throughout the whole course of updates, $ \DA $ iteratively utilizes two queries: $ \ProbAgree$ (\cite[Algorithm 8]{CLMP2024Dynamic}) and $ \Heavy $ (\cite[Algorithm 9]{CLMP2024Dynamic}).
They show that each query answer correctly with probability at least $ 1-n/2 $ (\cite[Lemma 49, Lemma 50, Lemma 52, Lemma 53]{CLMP2024Dynamic}).

To obtain $ \ADA $, we replace each decision of the randomized tests ProbabilisticAgreement and Heavy by the majority of three independent executions.
One execution of Heavy includes its internal calls to ProbabilisticAgreement, which are left unchanged.
As shown in \cref{lem:da.amplified.tests}, amplification reduces the error probability of each test decision to at most $ 3n^{-6} $, with only constant-factor overhead per decision.
The resulting whole-run guarantee is stated below and proved in \cref{sub:amplify.da}.

\paragraph{Amplifying the randomized tests.}
Dynamic Agreement uses the randomized tests ProbabilisticAgreement and Heavy.
A response is called correct if it satisfies the acceptance and rejection guarantees established for that test in \cite[Appendix~E]{CLMP2024Dynamic}, with the corrections described above.
Either response is permitted in the intermediate regimes where these guarantees prescribe no answer.

\begin{lemma}[Amplified agreement and heaviness tests]
    \label{lem:da.amplified.tests}
    Replace each test decision by the majority of three independent executions of the corresponding test.
    Conditional on the history before the decision, the amplified test gives a correct response with probability at least $ 1-3n^{-6} $.
    Amplification incurs only a constant-factor overhead per test decision.
\end{lemma}

\begin{proof}
    By \cite[Lemmas~49 and~50]{CLMP2024Dynamic}, one execution of ProbabilisticAgreement gives an incorrect response with probability at most $ n^{-3} $.
    The analogous guarantee for Heavy follows from \cite[Lemmas~52 and~53]{CLMP2024Dynamic}, with the corrected rejection guarantee.

    Conditioning on the preceding history fixes the graph and the test input.
    Since the three executions use independent fresh randomness, their errors are independent under this conditioning.
    The majority response can be incorrect only if at least two executions are incorrect.
    A union bound over the three pairs bounds this probability by $ 3n^{-6} $.

    Each amplified decision performs three executions of the original test.
    One execution of Heavy includes its internal calls to ProbabilisticAgreement, which are left unchanged.
\end{proof}

\begin{lemma}[Number of test invocations,
    {\cite[proof of Theorem~41, Appendix~C]{CLMP2024Dynamic}}]
    \label{lem:da.test.count}
    Fix an oblivious insertion-only stream of length at most $ n $.
    Let $ Q_t $ denote the number of agreement and heaviness test invocations during update $ t $.
    Then $ \mathbb E[Q_t]=O(\polylog(n)) $.
    Consequently, for every prefix of length $ t $,
    \begin{align}
        \mathbb E\left[\sum_{s=1}^{t}Q_s\right]
        =O(t\polylog(n)).
    \end{align}
\end{lemma}

\paragraph{The good event for Dynamic Agreement.}
Fix an oblivious insertion-only stream of length $ m\le n $.
Following the proof of \cite[Theorem~6]{CLMP2024Dynamic}, we separate test correctness from the remaining structural events.

Let $ \mathcal E_{\mathrm{test}}(t) $ denote the event that every amplified test decision through update $ t $ gives a correct response.
For each fixed prefix $ t\le m $, we use the following two structural events from that analysis:
\begin{enumerate}
    \item $ \mathcal A_t $ is the event that every reference cluster considered in \cite[Theorem~23]{CLMP2024Dynamic} is contained in one component of the maintained sparse graph.
    The bound in that theorem, together with the small-cluster correction and a union bound over at most $ n $ reference clusters, gives
    \begin{align}
        \Pr[
            \mathcal E_{\mathrm{test}}(t)
            \cap \overline{\mathcal A_t}
        ]
        \le n^{-7}.
    \end{align}

    \item $ \mathcal R_t $ is the aggregate regularity event used in the density analysis of \cite[Definition~28 and Observation~29]{CLMP2024Dynamic}.
    The estimates in \cite[Lemmas~25 and~27]{CLMP2024Dynamic}, aggregated over the relevant vertices and indices as in that analysis, give
    \begin{align}
        \Pr[\overline{\mathcal R_t}]
        \le n^{-100}.
    \end{align}
\end{enumerate}

On $ \mathcal E_{\mathrm{test}}(t)\cap\mathcal A_t\cap\mathcal R_t $, the argument in \cite[proof of Theorem~6]{CLMP2024Dynamic} implies that the sparse connected components give an $ \edFactor $-approximation, for an absolute constant $ \edFactor $.
Moreover, on $ \mathcal E_{\mathrm{test}}(t)\cap\mathcal R_t $, the density analysis in \cite[proof of Theorem~39]{CLMP2024Dynamic} gives the anchor-density property stated in \cref{lem:ccclass.is.dense}.

Define
\begin{align}
    \successDA
    :=
    \mathcal E_{\mathrm{test}}(m)
    \cap
    \bigcap_{t=1}^{m}
    (\mathcal A_t\cap\mathcal R_t).
\end{align}

\begin{lemma}[Anchor density, adapted from {\cite[proof of Theorem~39]{CLMP2024Dynamic}}]
    \label{lem:ccclass.is.dense}
    On $ \successDA $, for every time $ t\le m $, every component $ \ccclass $ of $ \sparseG{t} $ with $ \card{\ccclass}>1 $, and every $ a\in \ccclass\cap\anchorSet{t} $,
    \begin{align}
        |\openN{\sparseG{t}}{a}|\ge\tfrac34\card{\ccclass}-1.
    \end{align}
\end{lemma}

\begin{restatable}{lemma}{successDAProb}
    \label{lem:da.whole.run.success}
    For sufficiently large $ n $, the amplified algorithm satisfies
    \begin{align}
        \Pr[\successDA]\ge1-n^{-2}.
    \end{align}
    On this event, its sparse connected components give an $ \edFactor $-approximation and satisfy the anchor-density property simultaneously at every prefix.
\end{restatable}

\begin{proof}
    The insertion-only algorithm makes at most $ B_n=O(n^2\log^5 n) $ test decisions over the entire stream.
    This deterministic bound follows by counting test invocations in the algorithm and using \cite[Lemma~40]{CLMP2024Dynamic}.
    It holds independently of the test responses.

    By \cref{lem:da.amplified.tests} and a union bound over invocation indices,
    \begin{align}
        \Pr[\overline{\mathcal E_{\mathrm{test}}(m)}]
        \le 3B_n n^{-6}.
    \end{align}
    Combining this with the imported structural bounds gives
    \begin{align}
        \Pr[\overline{\successDA}]
        &\le
        \Pr[\overline{\mathcal E_{\mathrm{test}}(m)}]
        +
        \sum_{t=1}^{m}
        \Pr[
            \mathcal E_{\mathrm{test}}(t)
            \cap\overline{\mathcal A_t}
        ]
        +
        \sum_{t=1}^{m}
        \Pr[\overline{\mathcal R_t}] \\
        &\le 3B_n n^{-6}+n^{-6}+n^{-99} \\
        &=O(n^{-4}\log^5 n)+n^{-6}+n^{-99}
        \le n^{-2}
    \end{align}
    for sufficiently large $ n $.
    The approximation and density conclusions follow from the implications stated above.
\end{proof}

% subsection Auxiliary Lemmas From~ (end)


\subsection{Description of \texorpdfstring{$ \LabelSync $}{LabelSync} Algorithm}\label{sub:description_of_LabelSync} % (fold)

% \paragraph{\texorpdfstring{$ \LabelSync $}{LabelSync} Algorithm}\label{par:_texorpdfstring_labelsync_labelsync_algorithm} % (fold)
At the end of update $ t $, $ \LabelSync $ maintains a partition $ \ccclasses{t} $ of the current vertex set $ V_t $.
We call each member of $ \ccclasses{t} $ a \emph{class}.
Each class $ \ccclass $ has a distinct label $ \complabel(\ccclass) $, and every vertex $ w\in\ccclass $ stores this label as $ \complabel_t(w)=\complabel(\ccclass) $.
The partition $ \ccclasses{t} $ is the solution reported by $ \explicitDA $.
At the end of each update $ t $, the algorithm reports only the vertices that changed clusters, that is, $ (v_t,\complabel_t(v_t)) $ and every $ (w,\complabel_t(w)) $ with $ w\in V_{t-1} $ and $ \complabel_t(w)\ne\complabel_{t-1}(w) $.
We prove in \cref{sub:analysis_for_labelsync} that, on the success event $ \successDA $, the classes in $ \ccclasses{t} $ are exactly the connected components of $ \sparseG{t} $.
% paragraph LabelSync Algorithm (end)

\tienlong{Brevity: $ \ADA(t) $}
\paragraph{Overview of \texorpdfstring{$ \LabelSync $}{LabelSync}}\label{par:overview.of.LabelSync} % (fold)
The input to $ \LabelSync $ includes the sparse graph $ \sparseG{t} $ (after the update from time $ t-1 $ to time $ t $), the journal $ \journal $ of update $ t $, the arriving vertex $ v_t $, the data structure of $ \LabelSync $ at time $ t-1 $.

The goal is to update the data structure to reflect the connected components of $ \sparseG{t} $ and to report the vertices whose labels changed.

On a high level, the algorithm achieves this by performing a \emph{split} and a \emph{merge} operation.
Suppose that at time step $ t-1 $, $ \LabelSync $ successfully maintained the $ \ccclasses{t-1} $ as the connected components of $ \sparseG{t-1} $.
At time step $ t $, vertices in some components in $ \ccclasses{t-1} $ may not belong to the same components in $ \sparseG{t} $, and some vertices in different components in $ \ccclasses{t-1} $ may belong to the same component in $ \sparseG{t} $.
To fix this, we first perform the split operation, which partitions each old class $ \ccclass \in \ccclasses{t-1} $ into \piece{s}, where each \piece{} is a subset of vertices of $ \ccclass $ that belong to the same component in $ \sparseG{t} $.
After that, the merge operation groups \piece{s} that belong to the same component in $ \sparseG{t} $ together.
Finally, we update the labels of vertices in the new classes and report the vertices whose labels changed.

A direct approach to perform the split and merge operations would traverse the entire graph $ \sparseG{t} $ to find the connected components, which in the worst case takes $ O(n\polylog(n)) $ time.
To avoid this, we follow vertices and edges affected by $ \ADA(t) $ to perform the split and merge operations.
Note that, the length of the journal $ \journal{t} $ is at most the update time of $ \ADA $, therefore, we can afford to scan the whole journal for each update $ t $.
By reading off $ \journal{t} $, the algorithm can identify edges and vertices that were updated after the execution of $ \ADA $ for update $ t $.
It suffices to focus on the edges that were deleted to perform the split operation.
In particular, consider a component $ \ccclass \in \ccclasses{t-1} $.
If there are many edges inside $ \ccclass $ deleted by $ \ADA(t) $, we can safely scan the entire component $ \ccclass $ to split it into \piece{s}.
On the other hand, if only a few edges inside $ \ccclass $ were deleted, we demonstrate that in this case, all the vertices in $ \ccclass \cap \anchorSet{t-1} $ remain in the same component in $ \sparseG{t} $.
With the invariant~\eqref{eq:invariant.anchor} being maintained, one endpoint of each deleted edge is in $ \anchorSet{t-1} $.
This enables us to perform the split operation for $ \ccclass $ solely on the deleted edges and their endpoints not in $ \anchorSet{t-1} $.
After the split step, the merge step now operates on the \piece{s} created in the split step and the other components in $ \ccclasses{t-1} $ that have edges added by $ \ADA(t) $.
While merging \piece{s} in the same component of $ \sparseG{t} $ together, we take the label of a largest \piece{} in the component and assign it to all the vertices in the component.
Under this rule, intuitively, every time a vertex changes its label, the size of its component at least doubles.
Thus, on avarage, each vertex can change its label at most $ O(\log n) $ times, which gives a bound on the number of vertices reported by $ \LabelSync $.

% paragraph Overview of LabelSync (end)

\paragraph{Detailed Description of \texorpdfstring{$ \LabelSync $}{LabelSync}}\label{par:detailed_description_of_labelsync_} % (fold)
We provide the algorithm in more details with its main steps, data structure and the implementation.
The summary of the algorithm is given in \cref{alg.label.sync}.

Prior to the execution of $ \LabelSync{t} $, each class $ \ccclass \in \ccclasses{t-1} $ is represented by a record containing its label, size, and a doubly linked list of its vertices.
A set data structure is maintained to store the records of classes in $ \ccclasses{t-1} $.
We refer to this set as $ \setClasses$.
Each vertex $ w $ of $ \ccclass $ keeps a record of its label, a pointer to the record of $ \ccclass $, a pointer to $ w $'s position in $ \ccclass $'s doubly linked list, $ w $'s membership in $ \anchorSet{t-1} $, and its degree $ \deg_{t-1}(w):=\card{\openN{\sparseG{t-1}}{w}} $ in the sparse graph.

The algorithm consists of the following four steps.

\noindent\setword{\textbf{Step 1 (Identify Updated Edges)}}{StepOneLabelSync}:
In this step, we identify the edges that are updated after executing $ \ADA $ to handle update $ t $.
In particular, we compute
\begin{align}
    \delt{t}:=E(\sparseG{t-1})\setminus E(\sparseG{t}), \nonumber
    \qquad
    \addt{t}:=E(\sparseG{t})\setminus E(\sparseG{t-1}), \nonumber
\end{align}
the sets of edges deleted and added between $ \sparseG{t-1} $ and $ \sparseG{t} $ by $ \ADA $, respectively.
Note that, during the execution of $ \ADA $, an edge between two vertices $ u $ and $ v $ may be deleted and then added back, or added and then deleted, multiple times.
Here, $ \delt{t} $ and $ \addt{t} $ only record the edges that are present in $ \sparseG{t-1} $ (before update $ t $) but not in $ \sparseG{t} $, and vice versa, after $ \ADA $ finishes.

Additionally, for each old class $ \ccclass\in\ccclasses{t-1} $, let $ D_{\ccclass}:=\{(x,y)\in\delt{t}:x,y\in\ccclass\} $ be its deleted edges, and let $ d_{\ccclass}:=\card{D_{\ccclass}} $.
For each endpoint $ w $ of an edge in $ D_{\ccclass} $, let $ \delIncEdge{w}:=\card{\{e\in D_{\ccclass}:w\in e\}} $.

\textit{Implementation and data structures.}
We store $ \delt{t} $ and $ \addt{t} $ as lists, with each edge occurring at most once.
For each edge considered during the execution of $ \ADA $ for update $ t $, the journal $ \journal $ records its status in $ \sparseG{t-1} $ and $ \sparseG{t} $.
Hence, one pass over the journal establishes $ \delt{t} $ and $ \addt{t} $.
For each old class $ \ccclass $ with $ D_{\ccclass}\ne\emptyset $, we have a dictionary mapping its label $ \complabel(\ccclass) $ to the pair $ (D_{\ccclass},d_{\ccclass}) $, where $ D_{\ccclass} $ is a list of deleted edges and $ d_{\ccclass}=\card{D_{\ccclass}} $.
A second dictionary is employed to map each endpoint $ w $ of a deleted edge to its count $ \delIncEdge{w} $.
We fill both dictionaries in one pass over $ \delt{t} $.

\noindent\setword{\textbf{Step 2 (Split)}}{StepTwoLabelSync}:
We partition each old class into temporary sets called \emph{\piece{s}}, such that each \piece{} lies entirely within one component of $ \sparseG{t} $.

To do that, observe that conditioned on $ \successDA $, the \piece{s} of an old class $ \ccclass $ will be exactly
\begin{align}
    \mathcal P_{\ccclass}
    =\{\ccclass\cap K:K\text{ is a component of }\sparseG{t},\ \ccclass\cap K\ne\emptyset\}. \nonumber
\end{align}
Among the \piece{s} in $ \mathcal P_{\ccclass} $, we choose one to inherit the old label $ \complabel(\ccclass) $.
We refer to it as the \emph{retained \piece{}} of $ \ccclass $ and denote it by $ \retainedPiece{\ccclass} $.
The other \piece{s} initially have no assigned label.

If $ d_{\ccclass}=0 $, the class is represented by the single retained \piece{} $ \retainedPiece{\ccclass}=\ccclass $.
For each class with $ d_{\ccclass}>0 $, we distinguish two cases.
\begin{itemize}
    \item If $ d_{\ccclass}\ge\card{\ccclass}/2-1 $, we traverse $ \ccclass $ and group its vertices into \piece{s} such that, $ x,y \in \ccclass $ are grouped into the same \piece{} if and only if $ \comp{x}=\comp{y} $.

    A largest \piece{} is chosen as $ \retainedPiece{\ccclass} $.
    % Since $ \card{\ccclass}\le 2d_{\ccclass}+2\le 4d_{\ccclass} $, the number of vertices examined is proportional to $ d_{\ccclass} $.
    \item If $ d_{\ccclass}<\card{\ccclass}/2-1 $, choose an endpoint $ r\in\ccclass\cap\anchorSet{t-1} $ of an edge in $ D_{\ccclass} $.
    Such an endpoint exists because every old edge has an endpoint in $ \anchorSet{t-1} $, by definition of the anchor set.
    Define
    \begin{align}
        W_{\ccclass}:=\{w|\ \exists e\in D_{\ccclass}: w\in e, \ \text{and }
        w \in \ccclass\setminus\anchorSet{t-1},\ \delIncEdge{w}=\deg_{t-1}(w)\}, \nonumber
    \end{align}
    the set of vertices in $ \ccclass\setminus \anchorSet{t-1} $ whose incident edges at time $ t-1 $ are all deleted during update $ t $. \tienlong{Need to explain about $ \delIncEdge{w}=\deg_{t-1}(w) $}
    We will show that (\cref{lem:labelsync.robust.components}), on $ \successDA $, every vertex of $ \ccclass\setminus W_{\ccclass} $ remains connected to $ r $ after the deletions.
    For each vertex $ w\in W_{\ccclass} $, we group it with the same \piece{} as $ r $ if and only if $ \comp{w}=\comp{r} $.
    For the remaining vertices in $ W_{\ccclass} $, we group them into \piece{s}
    \begin{align}
        P_{\ccclass,\kappa}:=\{w\in W_{\ccclass}:\comp{w}=\kappa\}, \nonumber
    \end{align}
    for each $ \kappa\ne\comp{r} $ attained in $ W_{\ccclass} $.
    The retained \piece{} is
    \[
        \retainedPiece{\ccclass}:=\ccclass\setminus\bigcup_{\kappa\ne\comp{r}}P_{\ccclass,\kappa},
    \]
    where the union ranges only over the identifiers attained in $ W_{\ccclass} $.
    % We store $ R_{\ccclass} $ implicitly, with representative $ r $ and size $ \card{\ccclass}-\sum_{\kappa\ne\comp{r}}\card{P_{\ccclass,\kappa}} $.
    % The set $ W_{\ccclass} $ is found from the endpoints of $ D_{\ccclass} $, and its vertices are grouped using a hash map indexed by their component identifiers.
    % Thus this case does not require listing $ \ccclass $.
\end{itemize}

The \piece{s} of all old classes, together with the singleton $ \{v_t\} $, form a partition $ \mathcal P $ of $ V_t $. \tienlong{Not clear}

\textit{Implementation and data structures.}
For each class $ \ccclass $ participating in this step, we remove $ \ccclass $ from $ \setClasses $.

For each such $ \ccclass $, to group vertices of $ \ccclass $ with the same component identifier into one \piece{}, we utilize a dictionary mapping each identifier to the record containing a doubly linked list of vertices with that identifier.
We refer to this dictionary as $ \groupdict(\ccclass) $ for each such $ \ccclass $.

This dictionary is constructed as follows.

If $ d_{\ccclass}\ge\card{\ccclass}/2-1 $, the algorithm builds $ \groupdict(\ccclass) $ by traversing the vertices of $ \ccclass $.
In particular, for each vertex $ w\in \ccclass $, we add $ w $ to the doubly linked list of the record corresponding to the key $ \comp{w} $ in $ \groupdict(\ccclass) $.
We remove $ w $ from $ \ccclass $'s list, update $ w $'s pointers to the corresponding record in $ \groupdict(\ccclass) $ and its new position in the list of this record.
After the traversal, the record of a largest \piece{} in $ \groupdict(\ccclass) $ takes the label of $ \ccclass $, and the other records have no assigned label.

If $ 0<d_{\ccclass}<\card{\ccclass}/2-1 $, the algorithm first finds $ r $ and constructs the set $ W_{\ccclass} $ by examining each endpoints of $ D_{\ccclass} $.
We build $ \groupdict(\ccclass) $ just like in the previous case, but by traversing only the vertices in $ W_{\ccclass} $ instead of the whole class $ \ccclass $.
Also, during the traversal, we do not create a key in $ \groupdict(\ccclass) $ and the corresponding record for vertices with the same component identifier as $ r $, and also not remove them from the list of $ \ccclass $'s record.
After the traversal, the list of the record of the original class $ \ccclass $ now contains only the vertices of $ \retainedPiece{\ccclass} $.
Hence, we can use the record of $ \ccclass $ to represent $ \retainedPiece{\ccclass} $.
This record is added to $ \groupdict(\ccclass) $ with key $ \comp{r} $.

Note that, in both cases, whenever a vertex is moved from one list into another, we also update the size the corresponding records containing such lists.

% The singleton $ \{v_t\} $ initially has no assigned label.
% For each represented \piece{} $ P $, we store its size, a representative $ x_P\in P $, and its inherited label $ \lambda(P) $, where $ \lambda(P)=\bot $ if it has none.
% The inherited labels are temporary attributes of \piece{s}; vertex labels remain unchanged until Step 4.
% Parts other than implicitly represented retained \piece{s} are stored as explicit vertex lists.
% A hash map records the \piece{} containing each old vertex outside a retained part.
% Every other old vertex belongs to the retained \piece{} of its old class, so the \piece{} containing any given vertex can be found without listing its class.

\noindent\setword{\textbf{Step 3 (Merge)}}{StepThreeLabelSync}:
Let $ \piecesToMerge $ be a set containing all \piece{s} constructed for classes with $ d_{\ccclass}>0 $, the singleton $ \{v_t\} $, and the \piece{s} containing each endpoint of an edge in $ \addt{t} $.
This step groups \piece{s} in $ \piecesToMerge $ that belong to the same component of $ \sparseG{t} $ and assigns a label to each group.
Classes (\piece{s}) that are not in $ \piecesToMerge $ remain unchanged, and their labels are not modified.

For each $ \pieceToMerge \in \piecesToMerge $, define its \emph{component key} as $ \comp{x} $, for any $ x\in\pieceToMerge $.
The component key is well-defined because all vertices of a \piece{} share the same component identifier in $ \sparseG{t} $.

After all the \piece{s} in $ \piecesToMerge $ have been assigned a component key, we group the \piece{s} by their component keys.
For each such key $ \kappa $, choose a largest \piece{} $ \pieceToMerge^* $ with key $ \kappa $, breaking ties arbitrarily.
If $ \pieceToMerge^* $ has no assigned label, we assign it a previously unused label.
Vertices from the lists of all other \piece{s} with key $ \kappa $ are inserted into the list of $ \pieceToMerge^{*} $, which represents the new class.
The label of these vertices is updated to the label of $ \pieceToMerge^{*} $.

\textit{Implementation and data structures.}
We store $ \piecesToMerge $ in a set data structure containing records, initially containing a record for the singleton $ \{v_t\} $.
Other \piece{s} are added to $ \piecesToMerge $ as follows.

Recall that after~\ref{StepTwoLabelSync}, for each participating class $ \ccclass $, $ \groupdict(\ccclass) $ is a dictionary containing the records of all \piece{s} split from $ \ccclass $.
We add all such records to $ \piecesToMerge $ by traversing $ \groupdict(\ccclass) $.

For each endpoint $ x $ of an added edge (edges in $ \addt{t} $), based on its record, we can find the record of the \piece{} containing $ x $ and add it to $ \piecesToMerge $.
We also remove that record from $ \setClasses $, if it is currently in $ \setClasses $.

To group the \piece{s} in $ \piecesToMerge $ by their component keys, we use a dictionary mapping each key to a list of records of \piece{s} with that key.
We iterate over $ \piecesToMerge $, compute the component key of each \piece{}, and add its record to the list corresponding to that key in the dictionary.
We refer to this dictionary as $ \compdict $.

For each list of records with the same component key, we scan the list to find a largest \piece{} $ \pieceToMerge^* $ and assign it a new label if it does not have.
If $ \pieceToMerge^* $ is assigned a new label, so are its vertices.
After that, the lists of vertices of all the other \piece{s} in the same list with $ \pieceToMerge^* $ are inserted into the list of vertices of $ \pieceToMerge^* $, and the labels of these vertices are updated to the label of $ \pieceToMerge^* $.

Note that, records of all such vertices are also updated so that their pointers point to the record of the new class and the positions in the list.
The record of $ \pieceToMerge^{*} $ increases its size to reflect the number of vertices its list now contains.
$ \setClasses $ adds to itself the record of $ \pieceToMerge^* $.

To keep track of the vertices whose labels changed, we maintain a list $ \changedList $, which is initially empty.
While assigning labels to \piece{s} in the same component, we add a vertex with it label to $ \changedList $ whenever the \piece{} containing it is assigned a new label or when it is merged into another \piece{}.

\noindent\setword{\textbf{Step 4 (Report and Sync)}}{StepFourLabelSync}:
This step reports the newly arrived vertex $ v_{t} $ and its assigned label, and lists each vertex whose label changed together with its new label.
At the end, we also update the anchor memberships in $ \anchorSet{t} $ and degrees of vertices in $ \sparseG{t} $ for vertices with these updated information.

\textit{Implementation and data structures.}
Reporting the vertices with changed labels is proceeded by one pass over $ \changedList $, which keeps track of all vertices whose labels changed in~\ref{StepThreeLabelSync}.

To update the degrees in $ \sparseG{t} $ and membership in $ \anchorSet{t} $, we traverse the journal $ \journal $ and update the such values for only the vertices appearing in $ \journal $.
If $ v_t $ does not appear in $ \journal $, its degree is set to $ 0 $ and it is not a member of $ \anchorSet{t} $.

\tienlong{Remember to revise the input and output of the algorithm}
\begin{algorithm}[htbp]
  \small
  \DontPrintSemicolon
  \SetKwInOut{KwIn}{Input}
  \SetKwInOut{KwOut}{Output}
  \KwIn{The journal $ \journal $ of update $ t $, the arriving vertex $ v_t $, the stored state at time $ t-1 $, and connectivity queries in $ \sparseG{t} $.}
  \KwOut{$ (v_t,\complabel_t(v_t)) $ and every $ (w,\complabel_t(w)) $ with $ w\in V_{t-1} $ and $ \complabel_t(w)\ne\complabel_{t-1}(w) $.}

    $ \piecesToMerge\gets\{\{v_t\}\} $, $ \complabel(\{ v_{t} \})\gets \perp $

    \tcc{\ref{StepOneLabelSync}}
    \tcp{Make one scan over $ \journal $ to compute the following}
    $ \delt{t} \gets E(\sparseG{t-1})\setminus E(\sparseG{t}) $

    $ \addt{t} \gets E(\sparseG{t}) \setminus E(\sparseG{t-1}) $

    \tcp{Make one scan over $ \delt{t} $ to compute the following}
    $ D_{\ccclass} = \{(x,y)\in\delt{t}:x,y\in\ccclass\} $, $ d_{\ccclass} = \card{D_{\ccclass}} $ for each $ \ccclass \in \ccclasses{t-1} $ such that $ D_{\ccclass} \neq \emptyset $

    $ \delIncEdge{w}:=\card{\{e\in D_{\ccclass}:w\in e\}} $ for each $ w $ such that $ \delIncEdge{w} > 0 $

    \tcc{\ref{StepTwoLabelSync}}
    \ForEach{$ \ccclass $ with $ d_{\ccclass}>0 $}{\nllabel{line:labelsync.split.begin}
        \uIf{$ d_{\ccclass}\ge \card{\ccclass}/2-1 $}{\nllabel{line:labelsync.split.large.deletion.begin}
            \tcp{Make one scan over $ \ccclass $ to decompose it into \piece{s}.}
            $ \mathcal P_{\ccclass} \gets \{ \ccclass_{1}, \ccclass_{2}, \dots \} $ such that $ \ccclass = \cup_{\ccclass_{i}\in \mathcal P_{\ccclass}} \ccclass_{i} $, and $ x,y \in \ccclass_{i} $ if and only if $ \comp{x} = \comp{y} $


            $ \retainedPiece{\ccclass} \gets \argmax_{\ccclass_{i} \in \mathcal P_{\ccclass}}\card{\ccclass_{i}} $ is the retained \piece{} of $ \ccclass $ \tcp*[r]{Break ties arbitrarily.}
        \nllabel{line:labelsync.split.large.deletion.end}
        }
        \Else{\nllabel{line:labelsync.split.few.deletion.begin}
            \tcp{Make one scan over $ D_{\ccclass} $ to compute the following}
            One vertex $ r\in \ccclass\cap\anchorSet{t-1} $ incident to an edge of $ D_{\ccclass} $

            $ W_{\ccclass}\gets\{w\in\bigcup_{e\in D_{\ccclass}}e: w\notin\anchorSet{t-1},\ \delIncEdge{w}=\deg_{t-1}(w)\} $\;

            \tcp{Make one scan over $ W_{\ccclass} $ to decompose $ \ccclass $ into \piece{s}.}
            $ \mathcal P_{\ccclass} \gets \{ \ccclass_{1}, \ccclass_{2}, \dots \} $ such that $ \ccclass = \cup_{\ccclass_{i}\in \mathcal P_{\ccclass}} \ccclass_{i} $, and $ x,y \in \ccclass_{i} $ if and only if $ \comp{x} = \comp{y} $

            $ \retainedPiece{\ccclass} \gets \ccclass_{i} := \{ x: \comp{x} = \comp{r} \} $ is the retained \piece{} of $ \ccclass $
        }
        $ \complabel(\retainedPiece{\ccclass}) \gets \complabel(\ccclass) $

        $ \complabel(\ccclass_{i}) \gets \perp $ for all $ \ccclass_{i} \in \mathcal P_{ \ccclass } \setminus \{ \retainedPiece{\ccclass} \} $\;
        $ \piecesToMerge \gets \piecesToMerge \cup \{ \ccclass_{i}|\  \ccclass_{i}\in \mathcal P_{ \ccclass } \} $
    }\nllabel{line:labelsync.split.end}

    \tcc{\ref{StepThreeLabelSync}}
    \tcp{Make one scan over $ \addt{t} $ to add \piece{s} containing endpoints of edges in $ \addt{t} $ to $ \piecesToMerge $.}
    $ \piecesToMerge\gets\piecesToMerge\cup\{P: P \text{ is a \piece{} containing a vertex } x\in e\in \addt{t} \text{ and } P \text{ is not already in }\piecesToMerge\} $ \nllabel{line:labelsync.merge.begin} \nllabel{line:labelsync.piecesToMerge.include.addt}

    \BlankLine
    \tcp{Compute the component key of each \piece{} in $ \piecesToMerge $.}
    \ForEach{$ \pieceToMerge \in\piecesToMerge $}{ 
        $ \mathrm{key}(\pieceToMerge )\gets\comp{x} $, for any $ x\in\pieceToMerge $ \tcp*[r]{note that $ \comp{x} = \comp{y} $ for any $ x,y\in\pieceToMerge $} \nllabel{line:labelsync.assign.key}
    }

    $ \changedList \gets \{  \} $ \tcp*[r]{List to keep track of vertices whose labels changed.}
    \BlankLine
    \tcp{Group \piece{s} in $ \piecesToMerge $ by their component keys and assign labels.}
    \ForEach{distinct component key $ \kappa $ represented in $ \piecesToMerge $}{
        $ \piecesToMerge_\kappa\gets\{\pieceToMerge \in\piecesToMerge:\mathrm{key}(\pieceToMerge )=\kappa\} $\;\nllabel{line:labelsync.group.key}

        $ \pieceToMerge^* \gets \argmax_{\pieceToMerge \in\piecesToMerge_\kappa} \card{\pieceToMerge} $ \tcp*[r]{Break ties arbitrarily.}

        \uIf{$ \complabel(\pieceToMerge^*) = \perp $}{
            $ \complabel(\pieceToMerge^*) \gets \text{a previously unused label} $ \tcp*[r]{Assign a new label to $ \pieceToMerge^* $}

            Add $ (v, \complabel(\pieceToMerge^*)) $ to $ \changedList $, and update the label of $ v $ to $ \complabel(\pieceToMerge^*) $, for all $ v\in\pieceToMerge^* $
        }

        Add $ (v, \complabel(\pieceToMerge^*)) $ to $ \changedList $ for all $ v\in\pieceToMerge $, for all $ \pieceToMerge \in\piecesToMerge_\kappa\setminus\{\pieceToMerge^*\} $ \nllabel{line:labelsync.merge.add.small.pieces.changedlist}

        Merge $ \pieceToMerge $ to $ \pieceToMerge^* $ for all $ \pieceToMerge \in\piecesToMerge_\kappa\setminus\{\pieceToMerge^*\} $ and update the labels of vertices in $ \pieceToMerge $ to $ \complabel(\pieceToMerge^*) $


    }\nllabel{line:labelsync.merge.end}

    \tcc{\ref{StepFourLabelSync}}
    Make one scan over $ \journal{t} $ to update the degrees of vertices in $ \sparseG{t} $ and their membership in $ \anchorSet{t} $

    \Return{$ \changedList $}
  \caption{$ \LabelSync $ at the end of update $ t $.}
  \label{alg.label.sync}
\end{algorithm}
% paragraph Detailed Description of LabelSync (end)


% subsection Description of LabelSync (end)

\subsection{Analysis for \texorpdfstring{$ \LabelSync $}{LabelSync}}
\label{sub:analysis_for_labelsync}

\paragraph{Correctness.}
\label{par:correctness}
We use the invariant that every edge of $ \sparseG{t} $ has an endpoint in $ \anchorSet{t} $.
Let $ \successDA $ denote the success event from the analysis of $ \explicitDA $; on this event, the density bound in \cref{lem:ccclass.is.dense} holds at every time step.

\begin{lemma} \label{lem:labelsync.robust.components}
Let $ \sparseG $ be a simple undirected graph and let $ \anchorSet $ be a vertex set containing at least one endpoint of every edge of $ \sparseG $.
Let $ \ccclass $ be a connected component with $ \card{\ccclass}>1 $, and suppose that
\[
  |\openN{\sparseG}{a}|\ge\tfrac34\card{\ccclass}-1
  \qquad\text{for every }a\in \ccclass\cap\anchorSet.
\]
For any $ D\subseteq E(\sparseG[\ccclass]) $ with $ |D|<\card{\ccclass}/2-1 $, the graph $ \sparseG[\ccclass]\setminus D $ has one component containing $ \ccclass\cap\anchorSet $.
Every other component is a singleton outside $ \anchorSet $, and there are at most $ |D| $ such singletons.
\end{lemma}
\begin{proof}
    As $ \card{\ccclass} > 1 $, $ \ccclass $ has an least one edge.
    By definition, $ \anchorSet $ contains at least one vertex of all edges.
    Hence, $ \ccclass \cap \anchorSet \neq \emptyset $.
    We proceed by cases.

    $ \bullet $ If $ \card{\ccclass \cap \anchorSet} = 1 $, suppose $ \ccclass \cap \anchorSet = \{ a \} $.
    By definition of $ \anchorSet $, this implies that all edges of $ \ccclass $ are incident on $ a $, namely, $ \ccclass $ is a star graph with $ a $ as the center.
    Hence, deleting $ D $ results in a graph with one component containing $ a $ and all other components are singletons, and there are at most $ \card{D} $ of them.

    $ \bullet $ If $ \card{\ccclass \cap \anchorSet} > 1 $, consider two vertices $ a \neq b \in \ccclass \cap \anchorSet $. We show that the number of edge-disjoint paths of length at most $ 2 $ between $ a $ and $ b $ in $ \sparseG[\ccclass] $ is more than $ \card{D} $.

    Note that, for any two vertices $ w, w' \in \openN{\sparseG}{a} \cap \openN{\sparseG}{b} $, the paths $ (a,w,b) $ and $ (a,w',b) $ are of length $ 2 $ and edge-disjoint.
    If $ (a,b) $ is not an edge in $ \sparseG $, the number of such that is $ \card{\openN{\sparseG}{a} \cap \openN{\sparseG}{b}} \geq 2(\frac{3}{4}\card{\ccclass}-1) - (\card{\ccclass} - 2) \geq \frac{1}{2}\card{C} > \card{D} $.
    If $ (a,b) $ is an edge in $ \sparseG $, then $ \card{\openN{\sparseG}{a} \cap \openN{\sparseG}{b}} + \geq 2(\frac{3}{4}\card{\ccclass}-1) - \card{\ccclass} \geq \frac{1}{2}\card{C} - 2 $.
    The number of edge-disjoint paths of length at most $ 2 $ between $ a $ and $ b $, which also includes $ (a,b) $ in this case, is $ \card{\openN{\sparseG}{a} \cap \openN{\sparseG}{b}} + 1 \geq \frac{1}{2}\card{\ccclass} - 1 > \card{D} $.
    Hence, after deleting the edges in $ D $, there is still a path between $ a $ and $ b $.
    As this holds for any two vertices in $ \ccclass \cap \anchorSet $, all vertices in $ \ccclass \cap \anchorSet $ are in the same connected component of $ \sparseG[\ccclass]\setminus D $.

    After deleting $ D $, for any vertex in $ \ccclass \setminus \anchorSet $ that is not in the same component as $ \ccclass \cap \anchorSet $, all of its incident edges must be in $ D $, as the other endpoint of any such edge is in $ \ccclass \cap \anchorSet $.
    This also implies that all other components are singletons, and there are at most $ \card{D} $ of them.

\end{proof}
\begin{lemma}
\label{lem:labelsync.classes.are.connected.components}
Conditioned on $ \successDA $, after every update $ t $, the classes maintained by $ \LabelSync $ are exactly the connected components of $ \sparseG{t} $.
\end{lemma}
\begin{proof}
    We prove by induction on time step $ t $.

    At the base case $ t=1 $, the lemma holds as both $ \sparseG{1} $ and the classes maintained by $ \LabelSync{1} $ are both $ \{v_{1}\} $.

    For the inductive step, assume that the lemma holds for time step $ t-1 $, that is, every class maintained by $ \LabelSync{t-1} $ is a connected component of $ \sparseG{t-1} $.

    We show that it also holds for time step $ t $.

    $ \bullet $ First, we demonstrate that, after the split step (\cref{line:labelsync.split.begin}-\cref{line:labelsync.split.end}), every vertex in same \piece{} is in the same connected component of $ \sparseG{t} $. 

    Indeed, for each class $ \ccclass $ not participating in the split step ($ d_{\ccclass} = 0 $), they have no deleted edges.
    By induction hypothesis, $ \ccclass $ is a connected component of $ \sparseG{t-1} $. With none of its edges deleted, $ \ccclass $ remains a connected component of $ \sparseG{t} $.
    Hence, all its vertices are in the same connected component of $ \sparseG{t} $.

    It remains to consider classes $ \ccclass $ participating in the split step ($ d_{\ccclass} \geq 1 $).
    If $ d_{\ccclass} \geq \frac{1}{2}\card{\ccclass}-1 $, the algorithm iterates over all the vertices in $ \ccclass $, grouping vertices in the same connected component of $ \sparseG{t} $ together (\cref{line:labelsync.split.large.deletion.begin}-\cref{line:labelsync.split.large.deletion.end}).
    Note that, determining whether two vertices $ x,y $ are in the same connected component of $ \sparseG{t} $ can be done by checking if $ \comp{x} = \comp{y} $. 
    If $ d_{\ccclass} < \frac{1}{2}\card{\ccclass} - 1 $, conditioned on $ \successDA $, by~\cref{lem:ccclass.is.dense}, we have $ \openN{\sparseG{t-1}}{a} \geq \frac{3}{4}\card{\ccclass}-1 $ for any $ a \in \ccclass \cap \anchorSet{t-1} $.
    Now, by \cref{lem:labelsync.robust.components}, all vertices in $ \ccclass \cap \anchorSet{t-1} $ are in the same connected component of $ \sparseG{t} $, and all other vertices in $ \ccclass $ that are not in the same component as $ \ccclass \cap \anchorSet{t-1} $ are singletons.

    $ \bullet $ To finish the proof, it suffices to show that, the merge step (\cref{line:labelsync.merge.begin}-\cref{line:labelsync.merge.end}) merges all \piece{s} that are in the same connected component of $ \sparseG{t} $ and assigns them the same label, and assigns different labels to \piece{s} in different components. 
    This step begins with a set $ \piecesToMerge $, consisting of all the \piece{s} created in the split step, and the singleton $ \{v_t\} $; then adding to $ \piecesToMerge $ all the \piece{s} containing endpoints of an edge in $ \addt{t} $ (\cref{line:labelsync.merge.begin}-\cref{line:labelsync.piecesToMerge.include.addt}).
    All the \piece{s} participating in the merge step are the \piece{s} in $ \piecesToMerge $.

    If a \piece{} does not participate in the merge step, that is, is not included in the set $ \piecesToMerge $, then all of its vertices are not incident to any edge in $ \addt{t} $.
    By induction hypothesis, this \piece{} is a connected component of $ \sparseG{t-1} $.
    As no edges are added to this \piece{}, it remains a connected component of $ \sparseG{t} $ and keeps its label.

    For each \piece{} $ \pieceToMerge \in \piecesToMerge $, \cref{line:labelsync.assign.key} assigns a key $ \mathrm{key}(\pieceToMerge) = \comp{x} $, for some $ x \in \pieceToMerge $.
    We have shown above that all vertices in $ \pieceToMerge $ are in the same connected component of $ \sparseG{t} $, and hence, they have the same component identifier.
    It follows that, two \piece{s} share the same key if and only if they are in the same connected component of $ \sparseG{t} $.
    By grouping \piece{s} by their keys, the algorithm merges all \piece{s} which forms exactly one connected component of $ \sparseG{t} $ (\cref{line:labelsync.group.key}-\cref{line:labelsync.merge.end}).

    Finally, we argue that different components receive different labels. Indeed, each old label is inherited by exactly one retained \piece{}, and fresh labels are distinct and have never been used before.
    The classes after update $ t $ are therefore exactly the components of $ \sparseG{t} $.
\end{proof}

\begin{lemma}\label{lem:number.of.updated.vertices}
    For every $ t\leq n $, the total number of label changes of arrived vertices after executing $ \LabelSync $ (\cref{alg.label.sync}) for the first $ t $ updates is at most $ t\log_{2}n + 4(\log_{2}n + 3)\sum_{t'\leq t} \card{\delt{t'}} $.
\end{lemma}
\newcommand{\newpiecesize}[1]{\sigma_{#1}}
\newcommand{\smallpiecesize}[1]{\pi_{#1}}
\newcommand{\potential}[1]{\mathrm{Budget}(#1)}
\begin{proof}
    Fix any update time $ t'\in [1,t] $.
    We count the number of label changes after executing $ \LabelSync $ for update $ t' $, that is, the number of vertices with different labels between time $ t'-1 $ and $ t' $.

    At a high level, we upper bound this quantity by the number of vertices participating in the split and merge steps of $ \LabelSync $ for update $ t' $.

    Formally, consider the execution of $ \LabelSync $ for update $ t' $.
    Let $ \newpiecesize{t'} $ be the total size of the new \piece{s} created in the split step (\cref{line:labelsync.split.begin}-\cref{line:labelsync.split.end}). 
    In the merge step (\cref{line:labelsync.merge.begin}-\cref{line:labelsync.merge.end}), denote by $ \smallpiecesize{t'} $ be the total size of all the \piece{s} in $ \piecesToMerge $ that are not $ \pieceToMerge^* $.

    We first claim that, the number of vertices with changed labels is at most $ \newpiecesize{t'} + \smallpiecesize{t'} $. Indeed, the vertices with changed labels must belong to:
    \begin{itemize}
        \item Either one of the \piece{s} created in the split step, which also includes \piece{s} $ \piecesToMerge^{*} $ that requires a fresh label in the merge steps. The total size of these \piece{s} is at most $ \newpiecesize{t'} $.
        \item Or one of the \piece{s} in $ \piecesToMerge $ that are not $ \pieceToMerge^* $, which has total size at most $ \smallpiecesize{t'} $.
    \end{itemize}

    It suffices to demonstrate that $ \sum_{t'\leq t} \newpiecesize{t'} + \sum_{t'\leq t}\smallpiecesize{t'} \leq t \log_{2}n + 4(\log_{2} n + 3) \sum_{t'\leq t}\card{\delt{t'}} $.

    \paragraph{Bounding $ \newpiecesize{t'} $.}\label{par:bounding.newpiecesize} % (fold)
    In the split step, if $ d_{\ccclass} \geq \card{\ccclass}/2-1 $, the algorithm traverses all vertices in $ \ccclass $ and groups them into \piece{s} according to their identifiers. 
    The total size of the new \piece{s} created in this case is therefore at most $ \card{\ccclass} \leq 2d_{\ccclass}+2 \leq 4d_{\ccclass} $.
    In the remaining case, if $ d_{\ccclass} < \card{\ccclass}/2-1 $, the new \piece{s} are created grouping vertices in $ W_{\ccclass} $ according to their identifiers.
    The total size of the new \piece{s} created in this case is at most $ \card{W_{\ccclass}} $.
    By definition of $ W_{\ccclass} $ and the invariant~\eqref{eq:invariant.anchor}, we have $ \card{W_{\ccclass}} \leq d_{\ccclass} $.
    In either cases, the total size of the new \piece{s} created for class $ \ccclass $ is at most $ 4d_{\ccclass} $.
    Thus, the total size of the new \piece{s} created in the split step is
    \begin{align}
        \newpiecesize{t'} \leq \sum_{\ccclass \text{ participated in the split step}} 4d_{\ccclass} = 4\card{\delt{t'}}. \label{eq:bound.newpiecesize}
    \end{align}
    % paragraph Bounding newpiecesize (end)

    \paragraph{Bounding $ \smallpiecesize{t'}.$}\label{par:bounding_smallpiecesize_t_} % (fold)
    We show an upper bound for $ \sum_{t'\leq t} \smallpiecesize{t'} $ instead of bounding $ \smallpiecesize{t'} $ individually for each $ t' $.

    Consider the following scenario: throughout the whole course of the algorithm, for all time step $ t' $, there are no classes participating in the split step.
    Observe the merge step, a vertex $ v $ changes its label only if the \piece{} containing $ v $ joins a \piece{} with a larger size.
    This operation at least doubles the size of the \piece{} containing $ v $.
    Since the split opertation is never performed, it follows that any vertex can change its label at most $ \log_{2}n $ times.

    To account for the split step, we employ a potential function argument, which enables us to track the number of label changes allowed for each vertex
    Intuitively, every newly arrived vertex $ v $ is assigned a potential $ \potential{\{v\}} $ of $ \log_{2}n $, which corresponds to the maximum number of label changes it can undergo if the class containing $ v $ never participates in the split step.
    A split step performed on a class reduces the potential of each vertex in that class, while a merge step increases the potential of each vertex in the classes that get merged into a larger class.

    Formally, we define the potential function as follows. For a partition $ \ccclasses $ of a set $ V $, the potential of a class $ \ccclass\in\ccclasses $ is defined as
    \begin{align}
        \potential{\ccclass} = \card{\ccclass}\log_{2} \frac{n}{\card{\ccclass}}, \nonumber
    \end{align}
    and the potential of the whole partition is
    \begin{align}
        \potential{\ccclasses} = \sum_{\ccclass\in\ccclasses} \potential{\ccclass} = \sum_{\ccclass\in\ccclasses} \card{\ccclass}\log_{2} \frac{n}{\card{\ccclass}}. \nonumber
    \end{align}
    For each newly arrived vertex $ v $, assign a potential $ \potential{\{v\}} = \log_{2}n $.

    We analyze the change in potential of the partition maintained by $ \LabelSync $ at following three points in time: $ t'-1 $, after the split step of update $ t' $, and after the merge step of update $ t' $.
    At time $ t'-1 $, the partition is $ \potential{\ccclasses{t'-1}} $, which is the partition of classes maintained by $ \LabelSync $ at time $ t'-1 $, serving as the input to $ \LabelSync{t'} $. The partition after the merge step of update $ t' $ is the output of $ \LabelSync{t'} $. The partition after the split step of update $ t' $, denoted by $ \widetilde{\ccclasses{t'}} $, consists of $ v_{t'} $, the new \piece{s} created in the split step, and the \piece{s} that are not split.

    The split step transforming $ \ccclasses{t'-1} $ into $ \widetilde{\ccclasses{t'}} $ increases the potential. In particular, for each class $ \ccclass \in \ccclass{t'-1} $ split into $ \pieceToMerge_{1}, \pieceToMerge_{2}, \dots  $, with $ \pieceToMerge_{1} $ being the retained \piece{}, the potential increases by
    \begin{align}
        \sum_{i} \card{\pieceToMerge_{i}}\log_{2} \frac{n}{\card{\pieceToMerge_{i}}} - \card{\ccclass}\log_{2} \frac{n}{\card{\ccclass}} &= \sum_{i}\card{\pieceToMerge_{i}} \log_{2}\dfrac{\card{\ccclass}}{\card{\pieceToMerge_{i}}} \nonumber \\
                                                                                                                                                 &\leq \card{\pieceToMerge_{1}}\log_{2}\left(1+\dfrac{\card{\ccclass}-\card{\pieceToMerge_{1}}}{\card{\pieceToMerge_{1}}}\right) + (\card{\ccclass} - \card{\pieceToMerge_{1}})\log_{2}n  \nonumber \\
                                                                                                                                                 &\leq (\card{\ccclass} - \card{\pieceToMerge_{1}})(\log_{2}e + \log_{2}n). \nonumber
    \end{align}
    Note that, newly arrived vertex $ v_{t'} $ contributes $ \log_{2}n $ to the potential, and $ \sum_{\ccclass\in \ccclass{t'-1}}\card{\ccclass} - \card{\pieceToMerge_{1}} \leq \newpiecesize{t'} $.
    Hence,
    \begin{align}
        \potential{\widetilde{\ccclasses{t'}}} - \potential{\ccclasses{t'-1}} \leq \newpiecesize{t'}(\log_{2}e + \log_{2}n) + \log_{2}n. \label{eq.split.increases.potential}
    \end{align}

    The merge step transforming $ \widetilde{\ccclasses{t'}} $ into $ \ccclasses{t'} $ decreases the potential, and we argue that the decrease is at least $ \smallpiecesize{t'} $. 
    Indeed, consider a merge of $ \pieceToMerge_{1}, \pieceToMerge_{2}, \dots, $ with $ \pieceToMerge^* $ into $ \ccclass' $. Here, $ \card{\pieceToMerge^*} \geq \pieceToMerge_{i} $, and hence $ \card{\ccclass}\geq 2 \card{\pieceToMerge_{i}} $.
    Thus, for each such merge, the potential decreases by at least
    \begin{align}
        \sum_{i} \card{\pieceToMerge_{i}}\log_{2} \frac{n}{\card{\pieceToMerge_{i}}} - \card{\ccclass'}\log_{2} \frac{n}{\card{\ccclass'}} &\geq \sum_{i}\card{\pieceToMerge_{i}} \log_{2}\dfrac{\card{\ccclass'}}{\card{\pieceToMerge_{i}}} \geq \sum_{i} \card{\pieceToMerge_{i}}. \nonumber
    \end{align}
    By definition, $ \smallpiecesize{t'} $ covers the total size of all such $ \pieceToMerge_{i} $.
    It follows that $ \potential{\widetilde{\ccclasses{t'}}} - \potential{\ccclasses{t'}} \geq \smallpiecesize{t'} $.

    Taking the sum over all $ t'\leq t $, and note that $ \potential{\ccclasses{0}} = 0 \leq \potential{\ccclasses{t}} $, we have
    { \small
    \begin{align}
        \sum_{t'\leq t} \smallpiecesize{t'} &\leq (\potential{\widetilde{\ccclasses{1}}} - \potential{\ccclasses{1}}) + (\potential{\widetilde{\ccclasses{2}}} - \potential{\ccclasses{2}}) + \dots + (\potential{\widetilde{\ccclasses{t}}} - \potential{\ccclasses{t}}) \nonumber \\
                                            &\leq (\potential{\widetilde{\ccclasses{1}}} - \potential{\ccclasses{0}}) + (\potential{\widetilde{\ccclasses{2}}} - \potential{\ccclasses{1}}) + \dots + (\potential{\widetilde{\ccclasses{t}}} - \potential{\ccclasses{t-1}}) \nonumber \\
                                            &\leq \sum_{t'\leq t} \newpiecesize{t'}(\log_{2}e + \log_{2}n) + t\log n, \label{eq:bound.smallpiecesize}
    \end{align} 
    }
    where the last inequality follows from~\eqref{eq.split.increases.potential}.

    To finish the proof, combining~\eqref{eq:bound.newpiecesize} and~\eqref{eq:bound.smallpiecesize}, we obtain
    \begin{align}
        \sum_{t'\leq t} \newpiecesize{t'} + \sum_{t'\leq t}\smallpiecesize{t'} &\leq \sum_{t'\leq t} 4\card{\delt{t'}} + \sum_{t'\leq t} \newpiecesize{t'}(\log_{2}e + \log_{2}n) + t\log n \nonumber \\
                                                                                  &\leq 4(\log_{2}n + 3)\sum_{t'\leq t}\card{\delt{t'}} + t\log n. \nonumber
    \end{align}
    % paragraph Bounding $ \smallpiecesize{t'} $ (end)
\end{proof}

\begin{lemma}\label{lem:labelsync.cost}
    Suppose that after the update time $ t $, the length of the journal is $ \jnTime{t} $ and the number of vertices whose cluster label changes is $ \nChangedVers{t} $. 
    \tienlong{May be later remove this number of vertices whose cluster label changes, it seems confusing when bounding the running tim of the algorithm, or may be not, for example, suffix tree finding k matches}
    Then the running time of $ \LabelSync $ is $ O(\jnTime{t}\log n + \nChangedVers{t}) $.
\end{lemma}

\begin{proof}
    Write $ D= \card{\delt{t}} $ and $ A= \card{\addt{t}} $.
    Since the journal records all the edges that are removed or added, we have $ D+A\leq \jnTime{t} $.

    Scanning the journal.
    Identifying the deleted and added edges takes $ O(\jnTime{t}) $ dictionary operations.
    Scanning $\delt{t}$ then computes the counters and groups deleted edges by old class in $O(D)$ operations.
    Neither operation scans all vertices or all classes.

    Processing splits. There are at most $D$ processed classes, since each has $d_C\ge1$, and $\sum_C d_C\le D.$
    For a class in the many-deletions case, $|C|\le4d_C$, so enumerating and grouping its vertices takes $O(d_C\log n)$ time.
    In the few-deletions case, scanning its deleted edges identifies $W_C$ and $r$ in $O(d_C)$ operations, and querying and grouping the stranded vertices takes $O(d_C\log n)$ time.
    The retained piece remains in the old class list without enumerating it in the few-deletions case.
    Moving vertices into new piece lists takes time proportional to their number.

    Consequently, the entire split step takes $O(D\log n)$ time.
    Moreover, if $\sigma_t$ is the total number of vertices in new pieces, then $\sigma_t\le4D.$

    Bounding and constructing $\piecesToMerge$.
    There are at most $\sigma_t\le4D$ new pieces, since each is nonempty, and at most $D$ retained pieces, one per processed class.
    Added-edge endpoints introduce at most $2A$ further pieces, and $v_t$ contributes one singleton.
    Hence $|\piecesToMerge|\le4D+D+2A+1=O(\jnTime{t}+1)$.
    Using the piece lookup and deduplication described above, constructing this collection takes $O(\jnTime{t}+1)$ dictionary operations.

    Grouping and choosing labels.
    Querying one representative per piece and grouping by component identifier takes $O(|\piecesToMerge|\log n)$ time.
    Finding a largest piece in each group and deciding its final label takes $O(|\piecesToMerge|)$ further operations.

    Enumerating relabelled vertices.
    Every piece is enumerated directly from its linked list only when its vertices receive a different label.
    This includes a winning new piece that receives a fresh label; the newly inserted vertex contributes one additional operation.
    Thus all enumeration takes $O(\nChangedVers{t}+1)$ time.

    Reporting and bookkeeping.
    Applying the collected moves and reporting labels takes $O(\nChangedVers{t}+1)$ operations.
    Updating flags and degrees for journal vertices takes $O(\jnTime{t}+1)$ operations.

    Combining these bounds gives $O((\jnTime{t}+1)\log n+\nChangedVers{t})$.
    Separately, there are at most $\jnTime{t}$ journalled edge operations, each requiring $O(\log^2 n)$ amortized connectivity-update time.
\end{proof}

% section Explicit Dynamic Correlation Clustering Algorithm for Insertion-Only Stream (end)
